\documentclass[12pt,a4paper]{article}
\usepackage[utf8]{inputenc}
\usepackage[T1]{fontenc}
\usepackage[ngerman]{babel}
\usepackage[margin=2.5cm]{geometry}
\usepackage{tabularx}
\usepackage{booktabs}
\usepackage{url}

\setlength{\parindent}{0pt}
\setlength{\parskip}{6pt}

\newcommand{\vers}[2]{%
  \par\hangindent=1.6em\hangafter=1
  \textbf{#1}\enspace #2\par}

\begin{document}

\begin{center}
  {\LARGE\bfseries Hauskreis}\\[6pt]
  {\large Fichtenstr.~2, 01097 Dresden}\\[10pt]
  \begin{tabular}{ll}
    \textbf{Leitung:} & Stefan Helmert \\
    \textbf{Datum:}   & 08.10.2026 \\
    \textbf{Uhrzeit:} & 18:00 Uhr \\
  \end{tabular}
\end{center}

\vspace{1em}

\section*{Ablauf}

\begin{tabularx}{\textwidth}{@{}lX@{}}
  \toprule
  \textbf{Zeit} & \textbf{Programmpunkt} \\
  \midrule
  18:00 Uhr & Gebet \\
  18:05 Uhr & Bibelstudium (Teil 1 von 2: Markus 15,20--32; Fortsetzung beim nächsten Termin) \\
  19:00 Uhr & In 2 Gruppen: Wie geht es uns? Gebetsanliegen? \\
  20:00 Uhr & Aufräumen \\
  \bottomrule
\end{tabularx}

\vspace{1.5em}

\section*{Bibeltext: Markus 15,20--41}

\section*{Jesu Kreuzigung und Tod}

\subsection*{Kreuzigung}

\vers{20}{Und als sie ihn verspottet hatten, zogen sie ihm den Purpurmantel aus und zogen ihm seine Kleider an. Und sie führten ihn hinaus, dass sie ihn kreuzigten.}

\vers{21}{Und zwangen einen, der vorüberging, Simon von Kyrene, der vom Feld kam, den Vater des Alexander und des Rufus, dass er ihm das Kreuz trage.}
\vers{22}{Und sie brachten ihn zu der Stätte Golgatha, das heißt übersetzt: Schädelstätte.}
\vers{23}{Und sie gaben ihm Myrrhe im Wein zu trinken; aber er nahm's nicht.}
\vers{24}{Und sie kreuzigten ihn. Und sie teilten seine Kleider und warfen das Los darum, wer was bekommen sollte.}
\vers{25}{Und es war die dritte Stunde, als sie ihn kreuzigten.}
\vers{26}{Und es stand geschrieben, welche Schuld man ihm gab, nämlich: Der König der Juden.}
\vers{27--28}{Und sie kreuzigten mit ihm zwei Räuber, einen zu seiner Rechten und einen zu seiner Linken.}

\newpage

\subsection*{Verspottung am Kreuz}

\vers{29}{Und die vorübergingen, lästerten ihn und schüttelten ihre Köpfe und sprachen: Ha, der du den Tempel abbrichst und baust ihn auf in drei Tagen,}
\vers{30}{hilf dir nun selber und steig herab vom Kreuz!}
\vers{31}{Desgleichen verspotteten ihn auch die Hohenpriester untereinander samt den Schriftgelehrten und sprachen: Er hat andern geholfen und kann sich selber nicht helfen.}
\vers{32}{Der Christus, der König von Israel, er steige nun vom Kreuz, damit wir sehen und glauben. Und die mit ihm gekreuzigt waren, schmähten ihn auch.}

\subsection*{Finsternis, Schrei und Bekenntnis}

\vers{33}{Und zur sechsten Stunde kam eine Finsternis über das ganze Land bis zur neunten Stunde.}
\vers{34}{Und zu der neunten Stunde rief Jesus laut: Eli, Eli, lama asabtani? Das heißt übersetzt: Mein Gott, mein Gott, warum hast du mich verlassen?}
\vers{35}{Und einige, die dabeistanden, als sie das hörten, sagten: Siehe, er ruft den Elia.}
\vers{36}{Da lief einer, füllte einen Schwamm mit Essig, steckte ihn auf ein Rohr und gab ihm zu trinken, und sprach: Halt, lass sehen, ob Elia komme und ihn herabnehme.}
\vers{37}{Jesus aber rief laut und verschied.}
\vers{38}{Und der Vorhang im Tempel zeriss in zwei Stücke von oben bis unten.}
\vers{39}{Als aber der Hauptmann, der daneben stand, sah, dass er so verschieden war, sprach er: Wahrlich, dieser Mensch ist Gottes Sohn gewesen!}
\vers{40}{Es waren aber auch Frauen da, die von ferne zusahen, unter ihnen auch Maria Magdalena und Maria, die Mutter Jakobus' des Kleinen und des Joses, und Salome,}
\vers{41}{die ihm nachfolgten, als er in Galiläa war, und ihm dienten, und viele andere Frauen, die mit ihm nach Jerusalem hinaufgezogen waren.}

\vspace{1em}
{\footnotesize Bibeltext nach der Lutherbibel 2017, © Deutsche Bibelgesellschaft, Stuttgart.\\
Quelle: \url{https://www.bibleserver.com/LUT/Markus15}\\
Die Überschriften entsprechen exakt der Quellengliederung.\\
Ein Copy-Ausdruck für den Hauskreis ist gemäß Nutzungsbedingungen zulässig.}

\vspace{1.5em}

\newpage

\section*{Fragen zur Analyse}

\subsection*{Einstieg}

\textbf{Sprache}
\begin{enumerate}
  \item Welche Begriffe sind im Text unklar oder fremd? Sammelt sie und klärt sie mit Hilfe eines Bibellexikons (z.\,B. Bibel-Lexikon, apps.bibelwissenschaft.de).
\end{enumerate}

\textbf{Ort \& Zeit}
\begin{enumerate}\setcounter{enumi}{1}
  \item Warum heißt der Ort \glqq Schädelstätte\grqq{}, und warum liegt er außerhalb der Stadt (V. 22)?
  \item Sammle die Zeitangaben (V. 25. 33. 34) -- wie verläuft die Zeitleiste? Welches Fest ist im Hintergrund?
\end{enumerate}

\textbf{Kontext}
\begin{enumerate}\setcounter{enumi}{3}
  \item Was geschah unmittelbar vorher (Mk 14,43--15,19)?
  \item \glqq Sohn Gottes\grqq{}: Mk 1,1 $\leftrightarrow$ V. 39 -- wer erkennt es am Ende, und warum ist das überraschend?
\end{enumerate}

\textbf{Figuren}
\begin{enumerate}\setcounter{enumi}{5}
  \item Wer handelt, wer redet, wer schweigt in V. 20--41? (Wann spricht Jesus selbst?)
  \item Wer fehlt -- und warum? (Jünger, Mk 14,50!)
\end{enumerate}

\subsection*{Kreuzigung (V. 20--28)}
\begin{enumerate}\setcounter{enumi}{7}
  \item Simon wird \textbf{gezwungen}, das Kreuz zu tragen (V. 21) -- was heißt das für Mk 8,34 (\glqq sein Kreuz aufnehmen\grqq{})?
  \item Warum weist Jesus den Myrrhenwein ab (V. 23)?
  \item \glqq König der Juden\grqq\ (V. 26) -- als Anklage geschrieben: Inwiefern stimmt es?
\end{enumerate}

\subsection*{Verspottung am Kreuz (V. 29--32)}
\begin{enumerate}\setcounter{enumi}{10}
  \item Drei Gruppen spotten -- was wirft jede vor, wo wiederholt sich etwas?
  \item \textbf{Das Paradoxon von V. 31}: \glqq Er hat andern geholfen -- und kann sich selber nicht helfen.\grqq\ Wie kann beides zugleich wahr sein? Ist Nicht-sich-helfen-Können hier Schwäche -- oder gerade die Art, wie er anderen hilft?
  \item Wo wird der Spott unfreiwillig wahr (V. 29. 31. 32)?
  \item \glqq \ldots\ damit wir sehen und glauben\grqq\ (V. 32) -- warum ist das keine echte Glaubensbedingung?
\end{enumerate}

\subsection*{Finsternis, Schrei und Bekenntnis (V. 33--41)}
\begin{enumerate}\setcounter{enumi}{14}
  \item \glqq Warum hast du mich verlassen?\grqq\ (V. 34) -- nur Verzweiflung? (Ps 22: Klage, die ins Lob mündet.)
  \item Finsternis und Vorhangriss (V. 33. 38) -- wer handelt hier, was wird geöffnet bzw. gerichtet?
  \item Warum bekennt ausgerechnet der römische Hauptmann \glqq Gottes Sohn\grqq\ (V. 39)?
  \item Warum bleiben die Frauen und sehen alles (V. 40f) -- anders als die Jünger?
\end{enumerate}

\subsection*{Querbezüge zu anderen Bibelstellen}
\begin{enumerate}\setcounter{enumi}{18}
  \item Lies \textbf{Psalm 22} -- welche Verse unseres Textes klingen dort an? (Vergl. V. 29 mit Ps 22,8; V. 24 mit Ps 22,19; V. 34 mit Ps 22,2.)
  \item \textbf{Jesaja 53}: Welche Züge des \glqq leidenden Gottesknechtes\grqq\ (Jes 53,5.12) wiederholen sich in unserem Abschnitt?
  \item \textbf{Mk 8,31; 9,31; 10,33f}: Jesus kündigt sein Leiden dreimal an -- was davon wird in V. 20--27 buchstäblich erfüllt?
  \item \textbf{Mk 10,45}: \glqq \ldots\ sein Leben zu geben als Lösegeld für viele\grqq\ -- wie hört sich dieses Wort im Licht von V. 37.39 an?
  \item \textbf{Hebr 10,19f}: Was bedeutet der Vorhangriss (V. 38) im Licht dieses neutestamentlichen Deuteworts?
  \item \textbf{Amos 8,9 / Ex 10,21f}: Welches alte Motiv steht hinter der Finsternis (V. 33)?
\end{enumerate}

\subsection*{Abschluss}
\begin{enumerate}\setcounter{enumi}{24}
  \item Der Hauptmann erkennt Jesus im Sterbenden. Wo begegnen uns heute Menschen in \glqq gekreuzigten\grqq\ Situationen?
\end{enumerate}

\end{document}
